\chapter{验证与实验结果}
\label{chap:results}

本章报告本文加速器的正确性验证、板级端到端性能、优化消融、阶段计数器分析和资源时序结果。实验的重点不是将最终延迟包装为实时视频领先指标，而是验证本文提出的数据流和调度优化是否能够在真实 MPSoC 环境中降低 YOLOv3-tiny 链式推理的非计算开销，并解释性能瓶颈如何迁移。

\section{实验设置}

实验平台为 Xilinx Kria KV260，PL 器件为 XCK26，工具链为 Vivado/Vitis 2022.2。最终推荐硬件构建目录为 \texttt{D:/MPSoC/b\_kprefetch\_22}，核心时钟目标为 100 MHz。加速器参数为 ROWS=18、COLS=8、$Cout_{tile}=16$、IFM FIFO depth=1024，后端 raw-HWC full-tile cache 使用 URAM 承载 Conv5/6/8 的 $13\times13$ tile。

最终推荐软件配置为 \texttt{conv0\_conv9\_ddr\_demo}，开启 RawHwcConv4/5/6/8、early drain、pass prefetch、partial-PSUM overlap、continuous PSUM、backend full-tile 和 during-compute prefetch。验证路径包含两类：第一类是 batch-chain 模式，对 Conv0--Conv9 每层输出与 RTL semantic golden 做 bit-exact 比较；第二类是 DDR image demo，将 $416\times416\times3$ 图像包写入 DDR，由硬件执行十层推理并在 A53 侧进行 YOLO decode。

\section{正确性验证}

本文采用从模块到板级的多层验证。RTL 层面使用 Icarus Verilog 和 Vivado XSIM 覆盖 PE、array、feeder、AXI loader、PSUM buffer、collector、AXI-Lite bridge 和端到端 conv top。系统层面使用外部 golden 数据验证真实 Conv0、Layer06/Conv3、Conv4 以及后端 Conv5/6/8 raw-HWC tile。板级层面使用 KV260 batch-chain 和 DDR image demo 验证完整链路。

\begin{table}[htb]
\centering
\caption{关键正确性验证结果}
\label{tab:correctness}
\begin{tabular}{p{0.23\textwidth}p{0.49\textwidth}p{0.18\textwidth}}
\toprule
验证对象 & 覆盖内容 & 结果 \\
\midrule
Conv0 crop + pool & 真实输入 crop、卷积、LUT、pool、OFM packet & 0 mismatch \\
Layer06/conv3 pool & $52\times52\times64 \rightarrow 52\times52\times128$ & XSIM 与板级通过 \\
Conv4 pool & $26\times26\times128 \rightarrow 13\times13\times256$ & XSIM 与板级通过 \\
Conv5/6/8 raw-HWC & 后端 full-tile raw-HWC + overlap 配置 & XSIM 2726/0 \\
Conv0--Conv9 batch-chain & 十层链式 RTL semantic golden & PASS \\
DDR image demo & 动态图像输入与 YOLO decode & 两张图像 PASS \\
\bottomrule
\end{tabular}
\end{table}

需要强调的是，链式 golden 与单层 standalone golden 不可混用。由于前级硬件量化、激活和池化语义会影响后级输入，Conv5 之后的 expected output 必须由同一条 RTL semantic chain 生成。本文在 batch-chain 中逐层保存硬件输出并比较，从而验证层间数据重排、feature buffer 切换和检测头输出的一致性。

\section{端到端性能与优化消融}

表~\ref{tab:ablation} 汇总了主要板级优化版本的端到端 DDR demo 延迟。所有结果均来自 KV260 板级运行，输入图像和模型链路保持一致。不同版本的目标不是形成严格产品发布序列，而是展示本文优化如何逐步降低非计算开销。

\begin{table}[htb]
\centering
\caption{端到端 DDR demo 优化消融}
\label{tab:ablation}
\begin{tabular}{p{0.27\textwidth}p{0.43\textwidth}p{0.18\textwidth}}
\toprule
版本或优化阶段 & 主要变化 & 典型延迟 \\
\midrule
Batch/DDR baseline & batch stream，DDR 图像 demo 初版 & $\approx 1.18$ s \\
Weight 64-bit loader & 预打包权重与 64-bit PL 解包 & $\approx 0.861$ s \\
Drainpipe & 流水化 PSUM drain writer & $\approx 0.646$ s \\
Early drain + pass prefetch & 提前 drain 与下一 K pass 预取 & $\approx 0.387$ s \\
Continuous PSUM collector & 非 final pass 持续写 partial-PSUM & $\approx 0.371$ s \\
Backend full-tile & Conv5/6/8 使用 $13\times13$ full tile & $\approx 0.335$ s \\
During-compute prefetch & 当前 compute 中准备下一 K pass & $\approx 0.288$ s \\
\bottomrule
\end{tabular}
\end{table}

最终 \texttt{b\_kprefetch\_22} 构建完成 full-programming batch-chain 验证，日志中 UART detections 与 decode golden 匹配。两张 DDR image demo 的结果如表~\ref{tab:ddr_demo} 所示。

\begin{table}[htb]
\centering
\caption{最终 DDR image demo 结果}
\label{tab:ddr_demo}
\begin{tabular}{llll}
\toprule
图像 & 总延迟 & 检测结果 & 日志 \\
\midrule
maksssksksss0.png & 288.002 ms & with\_mask, score=0.357321 & 20260616\_193623 \\
maksssksksss1.png & 287.993 ms & with\_mask, score=0.295050 & 20260616\_193752 \\
\bottomrule
\end{tabular}
\end{table}

该结果约为 3.5 FPS，不应表述为通用实时检测领先性能。本文将其解释为在资源受限 XCK26 上，通过数据流与调度优化将早期约 1.18 s 的完整链式推理降低到约 0.29 s，获得约 4.1 倍端到端加速。

\section{阶段计数器分析}

仅报告端到端延迟无法说明优化来源。表~\ref{tab:stage_breakdown} 给出最终 DDR demo 的关键硬件计数器换算时间。可以看到，最终硬件 busy 时间约为 247.184 ms，端到端 total 与 PL busy 之间的差额来自 A53 runtime、DMA 管理、cache maintenance、OFM 重排和 YOLO decode 等软件侧开销。

\begin{table}[htb]
\centering
\caption{最终版本关键阶段计数器}
\label{tab:stage_breakdown}
\begin{tabular}{lr}
\toprule
指标 & 时间 \\
\midrule
Total DDR demo & 288.002 ms \\
HW busy & 247.184 ms \\
compute\_fire & 74.323 ms \\
feeder stage & 99.913 ms \\
compute stage & 121.022 ms \\
compute idle in stage & 46.699 ms \\
drain stage & 16.467 ms \\
\bottomrule
\end{tabular}
\end{table}

这些计数器说明，经过 drainpipe 和 continuous PSUM 后，drain 已不再是主要瓶颈；经过 backend full-tile 和 during-compute prefetch 后，剩余开销主要分布在 feeder/replay 和 compute stage idle。换言之，下一步优化不应继续单纯扩大阵列，而应继续提高 raw-HWC replay 吞吐、减少 compute stage 内部空泡，并评估 OFM debug packet 对软件重排的影响。

\section{资源与时序}

表~\ref{tab:resource} 给出最终推荐硬件构建的实现资源与时序。该版本在 Vivado 2022.2 下完成布局布线，WNS 为正，说明 100 MHz 目标可闭合。

\begin{table}[htb]
\centering
\caption{最终推荐构建资源与时序}
\label{tab:resource}
\begin{tabular}{lr}
\toprule
项目 & 数值 \\
\midrule
Build directory & \texttt{D:/MPSoC/b\_kprefetch\_22} \\
CLB LUTs & 84480 \\
CLB Registers & 53260 \\
BRAM Tile & 63 \\
URAM & 24 \\
DSP & 184 \\
WNS & +0.092 ns \\
TNS & 0 ns \\
WHS & +0.010 ns \\
Routing errors & 0 \\
\bottomrule
\end{tabular}
\end{table}

资源结果体现了本文优化的代价。raw-HWC full-tile cache 引入 URAM 以减少后端 tile 启动开销；continuous PSUM collector 和性能计数器增加 LUT/FF；during-compute prefetch 增加调度与队列控制逻辑。整体上，设计仍能装入 XCK26 并闭合 100 MHz，但 LUT 使用已经较高，继续增加 column-level PSUM 或更深 overlap 时需要谨慎评估资源和时序裕量。

\section{与相关工作的定性比较}

表~\ref{tab:qualitative_compare} 从研究重点角度比较本文与代表性工作。由于不同论文使用的模型、输入尺寸、器件、频率、量化策略和是否计入软件后处理差异较大，本文不进行不公平的绝对 FPS 横向排名。

\begin{table}[htb]
\centering
\caption{与代表性工作的定性比较}
\label{tab:qualitative_compare}
\begin{tabular}{p{0.18\textwidth}p{0.24\textwidth}p{0.22\textwidth}p{0.25\textwidth}}
\toprule
特征 & Angel-Eye~\cite{guo2019angeleye} & S2A~\cite{wei2017s2a} & 本文 \\
\midrule
主要目标 & 自动化映射流程 & systolic array 生成 & MPSoC 链式推理优化 \\
网络粒度 & 多模型映射 & 阵列综合 & YOLOv3-tiny Conv0--Conv9 \\
数据流讨论 & 编译映射为主 & 阵列结构为主 & DMA/IFM/PSUM/pass 调度 \\
板级端到端 & 有 & 非重点 & KV260 DDR image demo \\
性能解释 & 较少 & 较少 & stage/pass/counter 诊断 \\
\bottomrule
\end{tabular}
\end{table}

本文的优势在于系统闭环和瓶颈解释；不足在于通用性和绝对吞吐仍有限。当前 runtime 针对 YOLOv3-tiny 单尺度静态调度，不是通用 CNN 编译器；OFM 输出仍采用 debug packet 格式，验证友好但带宽效率不足；功耗数据尚未形成板级实测，后续可结合 Vivado power estimate 或外部功耗计进行能效评估。

\section{局限性与后续工作}

本文仍存在以下局限。第一，最终延迟约 288 ms，不能宣称实时视频级性能领先。其价值在于证明优化路径和系统完整性，而不是与 GPU/NPU 进行绝对性能竞争。第二，OFM debug packet 以每字节一个带地址 packet 的方式输出，便于验证但降低 DMA 带宽效率；后续应设计连续 HWC burst 写回。第三，raw-HWC cache 对层形状敏感，Conv3 实验说明并非所有层都适合启用该路径。第四，column-level PSUM 仍处于探索阶段，需要完整 top-level xsim 和板级 A/B 后才能作为正式优化。

后续工作可从四个方向展开：一是将 OFM 输出迁移为连续 HWC DMA burst，减少软件重排；二是继续优化 raw-HWC replay 吞吐和 compute stage idle；三是完善 column-level PSUM feedback，使 partial sum 按列粒度更早进入下一 pass；四是补充功耗实测和更广泛模型验证，提升论文结论的可比性和可推广性。

\section{本章小结}

本章给出了本文加速器的正确性、性能、资源和限制分析。实验表明，本文设计能够在 KV260/XCK26 上完成真实 Conv0--Conv9 链式推理，并通过 batch stream、raw-HWC cache、backend full-tile、continuous PSUM 和 during-compute prefetch 将端到端 DDR demo 延迟降低到约 0.29 s。性能计数器进一步说明，优化后瓶颈已从 PSUM drain 和权重加载转移到 feeder/replay 与 compute idle，为后续工作提供了明确方向。
